\chapter{Thiết kế hệ thống}

\section{Giới thiệu}

Sau khi hoàn thành quá trình phân tích yêu cầu, hệ thống AITasker được thiết kế
nhằm đảm bảo đáp ứng đầy đủ các yêu cầu chức năng và phi chức năng. Chương này
trình bày kiến trúc tổng thể, thiết kế các thành phần chính, mô hình lớp và cơ
sở dữ liệu của hệ thống.

%==============================================================

\section{Kiến trúc tổng thể}

AITasker được xây dựng theo mô hình Client--Server kết hợp kiến trúc nhiều tầng
(Multi-tier Architecture). Hệ thống gồm ba thành phần chính:

\begin{itemize}
\item Frontend: ReactJS.
\item Backend: FastAPI.
\item Database: PostgreSQL.
\end{itemize}

Ngoài ra, hệ thống còn tích hợp Google Gemini API để hỗ trợ các chức năng AI
như gợi ý mô tả dự án và chatbot hỗ trợ người dùng.

\begin{figure}[htbp]
\centering
\includegraphics[
width=\textwidth,
height=0.75\textheight,
keepaspectratio
]{architecture.png}
\caption{Kiến trúc tổng thể của hệ thống}
\label{fig:architecture}
\end{figure}

Frontend gửi yêu cầu đến Backend thông qua RESTful API. Backend xử lý nghiệp vụ,
thực hiện xác thực người dùng, thao tác với cơ sở dữ liệu và trả kết quả về giao
diện người dùng.

%==============================================================

\section{Thiết kế các thành phần}

Hệ thống được chia thành các module độc lập nhằm tăng khả năng mở rộng và bảo trì.

Các module chính gồm:

\begin{itemize}

\item Authentication

\item User Management

\item Enterprise

\item Expert

\item Project

\item Proposal

\item Contract

\item Payment

\item Review

\item Notification

\item Dashboard

\item AI Assistant

\end{itemize}

Mỗi module đảm nhiệm một nhóm chức năng riêng và giao tiếp với nhau thông qua
Service Layer.

%==============================================================

\section{Thiết kế lớp}

Hệ thống được xây dựng theo mô hình hướng đối tượng. Các lớp chính đại diện cho
các thực thể trong hệ thống như User, Enterprise, Expert, Project, Proposal,
Contract và Review.

\begin{figure}[htbp]
\centering
\includegraphics[
width=\textwidth,
height=0.80\textheight,
keepaspectratio
]{Class.png}
\caption{Biểu đồ lớp của hệ thống}
\label{fig:class}
\end{figure}

Biểu đồ lớp mô tả thuộc tính, phương thức và mối quan hệ giữa các đối tượng,
giúp việc phát triển và bảo trì hệ thống trở nên thuận lợi hơn.

%==============================================================

\section{Thiết kế cơ sở dữ liệu}

AITasker sử dụng PostgreSQL làm hệ quản trị cơ sở dữ liệu. SQLAlchemy ORM được
sử dụng để ánh xạ giữa các lớp trong chương trình và các bảng trong cơ sở dữ liệu.

Các bảng chính của hệ thống gồm:

\begin{itemize}

\item Users

\item Roles

\item Enterprises

\item Experts

\item Categories

\item Projects

\item Proposals

\item Contracts

\item Reviews

\item Notifications

\end{itemize}

\begin{figure}[htbp]
\centering
\includegraphics[
width=\textwidth,
height=0.80\textheight,
keepaspectratio
]{ERD.png}
\caption{Entity Relationship Diagram của hệ thống}
\label{fig:erd}
\end{figure}

ERD mô tả các thực thể, khóa chính, khóa ngoại và các mối quan hệ giữa các bảng,
đảm bảo dữ liệu được lưu trữ nhất quán và hạn chế dư thừa.

%==============================================================

\section{Thiết kế API và bảo mật}

Hệ thống sử dụng kiến trúc RESTful API để trao đổi dữ liệu giữa Frontend và
Backend. Dữ liệu được truyền dưới định dạng JSON.

Cơ chế bảo mật của hệ thống gồm:

\begin{itemize}

\item Xác thực bằng JWT.

\item Phân quyền theo vai trò (Administrator, Enterprise, Expert).

\item Mật khẩu được mã hóa bằng BCrypt.

\item Kiểm tra quyền truy cập tại các API.

\end{itemize}

Việc áp dụng JWT và phân quyền theo vai trò giúp đảm bảo an toàn và bảo mật
cho toàn bộ hệ thống.

%==============================================================

\section{Thiết kế tích hợp AI}

Hệ thống tích hợp Google Gemini API nhằm hỗ trợ người dùng trong quá trình sử
dụng hệ thống.

Các chức năng AI bao gồm:

\begin{itemize}

\item AI Job Assistant.

\item AI Chatbot.

\item Gợi ý mô tả dự án.

\item Hỗ trợ xây dựng yêu cầu dự án.

\end{itemize}

Việc tích hợp AI giúp nâng cao trải nghiệm người dùng và tăng hiệu quả trong
quá trình kết nối giữa doanh nghiệp và chuyên gia AI.

%==============================================================

\section{Kết luận}

Chương này đã trình bày thiết kế tổng thể của hệ thống AITasker, bao gồm kiến
trúc hệ thống, các module chức năng, biểu đồ lớp, mô hình cơ sở dữ liệu và cơ
chế bảo mật. Đây là cơ sở để triển khai hệ thống và kiểm thử trong các chương
tiếp theo.